\begin{proof}[Proof of \cref{obj:lem:adaptive-moment-contraction}]
All products over empty index sets equal one. We use the fixed-seed protocol experiment as the conditional-law version, including at seeds of zero \(\lambda\)-probability.

\begin{enumerate}
\item \textbf{Construction of the density.}
Define the observed-record and transcript spaces by
\[
\mathcal O_d=[d]\times\{0,1\}^2,
\qquad
\mathcal Z=\prod_{i=1}^n\mathcal Z_i.
\]
For \(o=(j_o,a_o,y_o)\in\mathcal O_d\), and for
\(\theta\in\mathbb R^d\) and
\(\mathbf o=(o_1,\ldots,o_n)\in\mathcal O_d^n\), set
\[
\mathsf{sgn}_{\mathrm{obs}}(o)
=(2a_o-1)(2y_o-1),
\qquad
\mathsf{mass}^{(1)}_\theta(o)
=\frac{1+\mathsf{sgn}_{\mathrm{obs}}(o)\theta_{j_o}}{4d},
\qquad
\mathsf{mass}^{(n)}_\theta(\mathbf o)
=\prod_{i=1}^n\mathsf{mass}^{(1)}_\theta(o_i).
\]
The reconstruction identities in \cref{obj:prop:signed-causal-bridge} give
\[
G_\theta(O=o)=\mathsf{mass}^{(1)}_\theta(o)
\qquad
\bigl(\theta\in[-1/2,1/2]^d\bigr).
\]
For every \(\mathbf o\in\mathcal O_d^n\) and \(r\in\mathcal R\), define the fixed-record transcript law by successive kernel integration:
\[
\mathsf{Law}^Q_{\mathbf o,r}(d\mathbf z)
=\prod_{i=1}^n Q_i(dz_i\mid o_i,\mathbf z_{<i},r).
\]
For \(\theta\in[-1/2,1/2]^d\), define the conditional transcript law and its zero-contrast reference by
\[
\mathsf{Prob}^Q_{\theta,r}
=\sum_{\mathbf o\in\mathcal O_d^n}
  \mathsf{mass}^{(n)}_\theta(\mathbf o)\,
  \mathsf{Law}^Q_{\mathbf o,r},
\qquad
\Lambda_r^Q=\mathsf{Prob}^Q_{0,r}.
\]
In particular,
\[
\Lambda_r^Q
=(4d)^{-n}\sum_{\mathbf o\in\mathcal O_d^n}
  \mathsf{Law}^Q_{\mathbf o,r}.
\]
Since \(d\ge2\), every coefficient in this last sum is positive. Consequently,
\[
\mathsf{Law}^Q_{\mathbf o,r}\ll\Lambda_r^Q
\qquad
\text{for every \(\mathbf o,r\)}.
\]

Apply the measurable-kernel-density hypothesis to these two probability-kernel families, parametrized by \((\mathbf o,r)\). It supplies jointly measurable functions
\[
\varphi^{\mathrm{ext}}_{\mathbf o}:
\mathcal R\times\mathcal Z\longrightarrow[0,\infty],
\qquad
\mathsf{Law}^Q_{\mathbf o,r}(E)
=\int_E\varphi^{\mathrm{ext}}_{\mathbf o}(r,\mathbf z)\,
  \Lambda_r^Q(d\mathbf z).
\]
Their integrals over \(\mathcal Z\) equal one, so their infinite-value sets are \(\Lambda_r^Q\)-null for each \(\mathbf o,r\). Define everywhere finite versions by
\[
\varphi_{\mathbf o}(r,\mathbf z)
=
\begin{cases}
\varphi^{\mathrm{ext}}_{\mathbf o}(r,\mathbf z),
 &\varphi^{\mathrm{ext}}_{\mathbf o}(r,\mathbf z)<\infty,\\
0,&\varphi^{\mathrm{ext}}_{\mathbf o}(r,\mathbf z)=\infty.
\end{cases}
\]
These versions are jointly measurable, nonnegative, and satisfy
\[
\mathsf{Law}^Q_{\mathbf o,r}(d\mathbf z)
=\varphi_{\mathbf o}(r,\mathbf z)\Lambda_r^Q(d\mathbf z),
\qquad
\int\varphi_{\mathbf o}(r,\mathbf z)\Lambda_r^Q(d\mathbf z)=1.
\]

Now define, on the full domain
\(\mathbb R^d\times\mathcal R\times\mathcal Z\),
\[
p^Q(\theta,r)(\mathbf z)
=
\sum_{\mathbf o\in\mathcal O_d^n}
 \mathsf{mass}^{(n)}_\theta(\mathbf o)\,
 \varphi_{\mathbf o}(r,\mathbf z).
\]
This is jointly measurable and real valued. On the parameter cube its summands are nonnegative, and finite summation of the density identities gives
\[
\mathsf{Prob}^Q_{\theta,r}(d\mathbf z)
=p^Q(\theta,r)(\mathbf z)\Lambda_r^Q(d\mathbf z).
\]

For \(j\in[d]\) and \(u\in\mathbb R\), define coordinate replacement by
\[
\theta^{[j:u]}_x
=
\begin{cases}
u,&x=j,\\
\theta_x,&x\ne j.
\end{cases}
\]
Each observed-record mass obeys
\[
\mathsf{mass}^{(1)}_{\theta^{[j:u]}}(o)
=
\mathsf{mass}^{(1)}_{\theta^{[j:0]}}(o)
+\frac{\mathbf1_{\{j_o=j\}}\mathsf{sgn}_{\mathrm{obs}}(o)}{4d}\,u.
\]
Thus each sample mass, and hence \(p^Q\), is a polynomial of degree at most \(n\) in each coordinate on all of \(\mathbb R\). Its coordinate coefficients are finite linear combinations of the integrable functions \(\varphi_{\mathbf o}(r,\cdot)\), so those coefficients and all coordinate derivatives are integrable against \(\Lambda_r^Q\).
% lean: CausalSmith.Stat.LdpOptvalueUniformFrontier.canonicalTranscriptDensity_integrable_coordinate_polynomial

\item \textbf{Pointwise normalized stage densities.}
Use the paired-input protocol \(\mathsf K\) supplied by
\cref{obj:prop:signed-causal-bridge}. It satisfies the same sequential privacy inequality at every history and seed. For
\(i\in[n]\), \(r\in\mathcal R\), and \(\eta\in\mathcal H_i\), define
\[
K_i^0(\,\cdot\mid r,\eta)
=\frac1{2d}\sum_{x=1}^d\sum_{s\in\{-1,1\}}
 \mathsf K_i(\,\cdot\mid(x,s),\eta,r).
\]
This is a probability kernel, and each summand is absolutely continuous with respect to it. The density hypothesis, with the same finite-value convention as above, therefore gives jointly measurable nonnegative real functions
\[
\psi^{\mathrm{raw}}_{i,x,s}:
\mathcal R\times\mathcal H_i\times\mathcal Z_i\longrightarrow[0,\infty)
\]
such that
\[
\mathsf K_i(d\zeta\mid(x,s),\eta,r)
=\psi^{\mathrm{raw}}_{i,x,s}(r,\eta,\zeta)
 K_i^0(d\zeta\mid r,\eta).
\]

For each \(r,\eta\), the definition of \(K_i^0\) and uniqueness of densities imply
\[
\sum_{x,s}\psi^{\mathrm{raw}}_{i,x,s}(r,\eta,\zeta)=2d
\quad
K_i^0(\cdot\mid r,\eta)\text{-almost everywhere}.
\]
Privacy gives, for every measurable \(E\subseteq\mathcal Z_i\),
\[
\int_E\psi^{\mathrm{raw}}_{i,x,s}(r,\eta,\zeta)
 K_i^0(d\zeta\mid r,\eta)
\le
e^\varepsilon
\int_E\psi^{\mathrm{raw}}_{i,x',s'}(r,\eta,\zeta)
 K_i^0(d\zeta\mid r,\eta).
\]
The comparison of these integrals over all measurable sets implies the corresponding density inequality almost everywhere.

Define the measurable good set
\[
\begin{split}
\mathcal G_i
=\biggl\{(r,\eta,\zeta):\;&
 \sum_{x,s}\psi^{\mathrm{raw}}_{i,x,s}(r,\eta,\zeta)=2d,\\
&\psi^{\mathrm{raw}}_{i,x,s}(r,\eta,\zeta)
 \le e^\varepsilon\psi^{\mathrm{raw}}_{i,x',s'}(r,\eta,\zeta)
 \ \text{for all }x,s,x',s'\biggr\},
\end{split}
\]
and set, on their full domains,
\[
\psi_{i,x,s}(r,\eta,\zeta)
=
\begin{cases}
\psi^{\mathrm{raw}}_{i,x,s}(r,\eta,\zeta),
 &(r,\eta,\zeta)\in\mathcal G_i,\\
1,&(r,\eta,\zeta)\notin\mathcal G_i.
\end{cases}
\]
For each \(r,\eta\), the good set has full reference measure. Hence this repair preserves every represented kernel and gives, everywhere,
\[
\sum_{x,s}\psi_{i,x,s}=2d,
\qquad
\psi_{i,x,s}\le e^\varepsilon\psi_{i,x',s'}.
\]
Summing the second inequality first over the left-hand input and then over the right-hand input yields
\[
2d\le 2d\,e^\varepsilon\psi_{i,x,s},
\qquad
2d\,\psi_{i,x,s}\le e^\varepsilon\sum_{x',s'}\psi_{i,x',s'}
=2d\,e^\varepsilon.
\]
Therefore
\[
e^{-\varepsilon}\le\psi_{i,x,s}(r,\eta,\zeta)\le e^\varepsilon
\qquad\text{everywhere}.
\]
% lean: CausalSmith.Stat.LdpOptvalueUniformFrontier.stage_density_repaired_of_gate

\item \textbf{The stage product and coordinate identification.}
For all \(\theta\in\mathbb R^d\) and all declared stage arguments, define
\[
\mathsf{dens}_i(\theta,r,\eta,\zeta)
=
\frac1{2d}\sum_{x,s}(1+s\theta_x)\psi_{i,x,s}(r,\eta,\zeta),
\qquad
\mathsf{slope}_{ij}(r,\eta,\zeta)
=
\frac{\psi_{i,j,1}(r,\eta,\zeta)
      -\psi_{i,j,-1}(r,\eta,\zeta)}{2d}.
\]
Normalization gives the identities
\[
\mathsf{dens}_i(\theta,r,\eta,\zeta)
=1+\sum_{j=1}^d\theta_j\mathsf{slope}_{ij}(r,\eta,\zeta),
\]
\[
\mathsf{dens}_i(\theta^{[j:u]},r,\eta,\zeta)
=
\mathsf{dens}_i(\theta,r,\eta,\zeta)
+(u-\theta_j)\mathsf{slope}_{ij}(r,\eta,\zeta).
\]
Define the reference stage-chain law and the stage product by
\[
\Lambda_r^{Q,\mathrm{stage}}(d\mathbf z)
=\prod_{i=1}^n K_i^0(dz_i\mid r,\mathbf z_{<i}),
\qquad
\mathsf{prod}^Q(\theta,r,\mathbf z)
=\prod_{i=1}^n
 \mathsf{dens}_i(\theta,r,\mathbf z_{<i},z_i).
\]

For a fixed paired-input vector
\(\boldsymbol\chi=((j_i,s_i))_{i=1}^n
 \in([d]\times\{-1,1\})^n\), define
\[
\mathsf{Law}^{\mathsf K}_{\boldsymbol\chi,r}(d\mathbf z)
=\prod_{i=1}^n
 \mathsf K_i(dz_i\mid(j_i,s_i),\mathbf z_{<i},r).
\]
Successive integration of the stage density identities gives
\[
\mathsf{Law}^{\mathsf K}_{\boldsymbol\chi,r}(d\mathbf z)
=
\left\{\prod_{i=1}^n
 \psi_{i,j_i,s_i}(r,\mathbf z_{<i},z_i)\right\}
 \Lambda_r^{Q,\mathrm{stage}}(d\mathbf z).
\]
The fair ancillary-bit reconstruction in
\cref{obj:prop:signed-causal-bridge}, followed by expansion of the averaged kernels at successive stages, gives the fixed-seed identity
\[
\mathsf{Prob}^Q_{\theta,r}
=
\sum_{\boldsymbol\chi}
 \left\{\prod_{i=1}^nP_\theta(j_i,s_i)\right\}
 \mathsf{Law}^{\mathsf K}_{\boldsymbol\chi,r}
\qquad
\bigl(\theta\in[-1/2,1/2]^d\bigr).
\]
Indeed, each ancillary assignment vector has probability \(2^{-n}\), and its recovered observed-record vector is obtained coordinatewise by
\[
o_i
=\left(j_i,a_i,\frac{1+s_i(2a_i-1)}2\right),
\qquad
(a_1,\ldots,a_n)\in\{0,1\}^n.
\]
Finite distribution of sums over products now gives
\[
\begin{split}
&\sum_{\boldsymbol\chi}
 \prod_{i=1}^n
 \left\{
 \frac{1+s_i\theta_{j_i}}{2d}
 \psi_{i,j_i,s_i}(r,\mathbf z_{<i},z_i)
 \right\}\\
&\hspace{3em}
=\prod_{i=1}^n
 \left\{\frac1{2d}\sum_{x,s}
 (1+s\theta_x)\psi_{i,x,s}(r,\mathbf z_{<i},z_i)\right\}
=\mathsf{prod}^Q(\theta,r,\mathbf z).
\end{split}
\]
Consequently,
\[
\mathsf{Prob}^Q_{\theta,r}(d\mathbf z)
=\mathsf{prod}^Q(\theta,r,\mathbf z)
 \Lambda_r^{Q,\mathrm{stage}}(d\mathbf z).
\]
At \(\theta=0\), every stage density equals one. Thus
\[
\Lambda_r^{Q,\mathrm{stage}}=\Lambda_r^Q,
\qquad
\mathsf{Prob}^Q_{\theta,r}(d\mathbf z)
=\mathsf{prod}^Q(\theta,r,\mathbf z)\Lambda_r^Q(d\mathbf z).
\]

Fix \(\theta\in[-1/2,1/2]^d\), \(r\), and \(j\). Both
\[
u\longmapsto p^Q(\theta^{[j:u]},r)(\mathbf z),
\qquad
u\longmapsto\mathsf{prod}^Q(\theta^{[j:u]},r,\mathbf z)
\]
are polynomials of degree at most \(n\). Define distinct interpolation points by
\[
\upsilon_\iota=\frac1{\iota+2},
\qquad \iota\in\mathbb N.
\]
Every \(\theta^{[j:\upsilon_\iota]}\) lies in the cube. Since the two nonnegative functions represent the same law there, uniqueness of densities gives equality almost everywhere at each \(\upsilon_\iota\). Taking the countable intersection of these full-measure sets, the two polynomials agree at infinitely many distinct points. Polynomial uniqueness therefore gives
\[
\text{for \(\Lambda_r^Q\)-almost every \(\mathbf z\), for every \(u\in\mathbb R\),}\qquad
p^Q(\theta^{[j:u]},r)(\mathbf z)
=\mathsf{prod}^Q(\theta^{[j:u]},r,\mathbf z).
\]
This identity also identifies every order of coordinate differentiation on the same full-measure set.
% lean: CausalSmith.Stat.LdpOptvalueUniformFrontier.transcript_coordinate_ae_eq_stageProduct

\item \textbf{Scores and their moments.}
Fix \(\theta\in[-1/2,1/2]^d\), \(r\), and \(j\). For each stage define
\[
\psi_{i,\min}(r,\eta,\zeta)
=\min_{x,s}\psi_{i,x,s}(r,\eta,\zeta).
\]
The weights \((1+s\theta_x)/(2d)\) are nonnegative and sum to one. Privacy and the minimum therefore imply
\[
\mathsf{dens}_i(\theta,r,\eta,\zeta)
\ge\psi_{i,\min}(r,\eta,\zeta)>0,
\]
\[
\left|\psi_{i,j,1}(r,\eta,\zeta)
      -\psi_{i,j,-1}(r,\eta,\zeta)\right|
\le(e^\varepsilon-1)\psi_{i,\min}(r,\eta,\zeta)
\le(e^\varepsilon-1)\mathsf{dens}_i(\theta,r,\eta,\zeta).
\]
Hence
\[
\left|
\frac{\mathsf{slope}_{ij}(r,\eta,\zeta)}
     {\mathsf{dens}_i(\theta,r,\eta,\zeta)}
\right|
\le\beta.
\]

Under \(\mathsf{Prob}^Q_{\theta,r}\), define expectation, the history filtration, and the scores by
\[
\mathbb E^Q_{\theta,r}[\cdot]
=\int_{\mathcal Z}(\cdot)\,\mathsf{Prob}^Q_{\theta,r}(d\mathbf z),
\qquad
\mathcal F_v=\sigma(Z_1,\ldots,Z_{\min\{v,n\}}),
\quad v\in\mathbb N,
\]
\[
\mathsf{score}^{\theta,r,j}_i(\mathbf z)
=
\begin{cases}
\dfrac{\mathsf{slope}_{ij}(r,\mathbf z_{<i},z_i)}
      {\mathsf{dens}_i(\theta,r,\mathbf z_{<i},z_i)},
 &1\le i\le n,\\[6pt]
0,&i>n,
\end{cases}
\qquad i\in\{1,2,\ldots\}.
\]
The parameter-specific stage kernel is
\[
\mathsf{dens}_i(\theta,r,\eta,\zeta)
 K_i^0(d\zeta\mid r,\eta).
\]
Its successive kernel integrations give
\(\mathsf{Prob}^Q_{\theta,r}\), by the product identity above. Each input-row density integrates to one, so
\[
\int\mathsf{slope}_{ij}(r,\eta,\zeta)
 K_i^0(d\zeta\mid r,\eta)
=\frac{1-1}{2d}=0.
\]
Cancellation of the positive stage density in the score integral gives
\[
\mathbb E^Q_{\theta,r}
 [\mathsf{score}^{\theta,r,j}_i\mid\mathcal F_{i-1}]
=0,
\qquad
|\mathsf{score}^{\theta,r,j}_i|\le\beta.
\]
The zero extension satisfies these identities after the horizon as well.

For \(k,v\in\mathbb N\), define the elementary score polynomials recursively by
\[
\mathsf{elem}_{0,v}=1,
\qquad
\mathsf{elem}_{k,0}=0\quad(k\ge1),
\]
\[
\mathsf{elem}_{k,v+1}
=\mathsf{elem}_{k,v}
+\mathsf{score}^{\theta,r,j}_{v+1}\mathsf{elem}_{k-1,v}
\quad(k\ge1).
\]
Here \(\mathsf{elem}_{k,v}\) is a function of the transcript, with the fixed
\(\theta,r,j\) understood. The recursion gives
\(\mathsf{elem}_{k,v}=0\) when \(k>v\), and each prefix polynomial is
\(\mathcal F_v\)-measurable. Conditional mean zero thus implies
\[
\mathbb E^Q_{\theta,r}
 \left[
 \mathsf{elem}_{k,v}\,
 \mathsf{score}^{\theta,r,j}_{v+1}\,
 \mathsf{elem}_{k-1,v}
 \right]=0
\qquad(k\ge1).
\]
Squaring the recursion and using the score bound yields
\[
\mathbb E^Q_{\theta,r}[\mathsf{elem}_{k,v+1}^2]
\le
\mathbb E^Q_{\theta,r}[\mathsf{elem}_{k,v}^2]
+\beta^2\mathbb E^Q_{\theta,r}[\mathsf{elem}_{k-1,v}^2].
\]
Induction from the displayed initial values, using
\(\binom{v+1}{k}=\binom vk+\binom v{k-1}\), gives
\[
\mathbb E^Q_{\theta,r}[\mathsf{elem}_{k,v}^2]
\le\binom vk\beta^{2k}.
\]
This includes \(k=0\) and \(k>v\).

The nonnegative centered square gives
\[
\begin{split}
0
&\le
\mathbb E^Q_{\theta,r}
 \left[
 \left(
 |\mathsf{elem}_{k,v}|
 -\mathbb E^Q_{\theta,r}|\mathsf{elem}_{k,v}|
 \right)^2
 \right]\\
&=
\mathbb E^Q_{\theta,r}[\mathsf{elem}_{k,v}^2]
-\left(\mathbb E^Q_{\theta,r}|\mathsf{elem}_{k,v}|\right)^2.
\end{split}
\]
Since \(\beta\ge0\), it follows that
\[
\mathbb E^Q_{\theta,r}|\mathsf{elem}_{k,v}|
\le\sqrt{\binom vk}\,\beta^k.
\]
% lean: CausalSmith.Stat.LdpOptvalueUniformFrontier.originalTranscript_scoreElementary_firstMoment

\item \textbf{The derivative certificate.}
Multiplying the affine coordinate-shift identities recursively gives
\[
\mathsf{prod}^Q(\theta^{[j:\theta_j+u]},r,\mathbf z)
=
\mathsf{prod}^Q(\theta,r,\mathbf z)
\sum_{q=0}^n u^q\mathsf{elem}_{q,n}(\mathbf z),
\qquad u\in\mathbb R.
\]
The simultaneous coordinate identity from the preceding construction therefore yields
\[
\partial_{\theta_j}^{\,k}p^Q(\theta,r)(\mathbf z)
=
k!\,p^Q(\theta,r)(\mathbf z)\,
\mathsf{elem}_{k,n}(\mathbf z)
\quad
\Lambda_r^Q\text{-almost everywhere}.
\]
Changing measure by the nonnegative density \(p^Q\), and applying the first-moment estimate, gives
\[
\begin{split}
\int
\left|\partial_{\theta_j}^{\,k}p^Q(\theta,r)(\mathbf z)\right|
 \Lambda_r^Q(d\mathbf z)
&=
k!\,\mathbb E^Q_{\theta,r}|\mathsf{elem}_{k,n}|\\
&\le k!\sqrt{\binom nk}\,\beta^k,
\qquad 1\le k\le n.
\end{split}
\]
Finally, the coordinate polynomial representation of the finite-sum density gives, pointwise on its full domain,
\[
\partial_{\theta_j}^{\,k}p^Q(\theta,r)(\mathbf z)=0
\qquad
\bigl(\theta\in\mathbb R^d,\ r\in\mathcal R,\
\mathbf z\in\mathcal Z,\ j\in[d],\ k>n\bigr).
\]
Together with the construction and representation already proved, these are all the asserted density properties.
% lean: CausalSmith.Stat.LdpOptvalueUniformFrontier.canonical_density_certificate_rows_of_gate

\item \textbf{A single-coordinate moment bound.}
Fix \(0<\alpha\le1/2\), the two supported probability priors, and
\(1\le k\le n\) with the asserted matching moments. For
\(\theta\in[-1/2,1/2]^d\), \(r\in\mathcal R\), and \(j\in[d]\), define
\[
\mathsf{diff}^{\theta,r}_j(\mathbf z)
=
\int p^Q(\theta^{[j:u]},r)(\mathbf z)\,\nu_0(du)
-\int p^Q(\theta^{[j:u]},r)(\mathbf z)\,\nu_1(du),
\]
\[
\mathsf{sign}^{\theta,r}_j(\mathbf z)
=
\begin{cases}
\dfrac{\mathsf{diff}^{\theta,r}_j(\mathbf z)}
      {|\mathsf{diff}^{\theta,r}_j(\mathbf z)|},
 &\mathsf{diff}^{\theta,r}_j(\mathbf z)\ne0,\\[6pt]
0,&\mathsf{diff}^{\theta,r}_j(\mathbf z)=0,
\end{cases}
\]
and the scalar polynomial, defined for every \(u\in\mathbb R\), by
\[
\mathsf{poly}^{\theta,r}_j(u)
=
\int
 \mathsf{sign}^{\theta,r}_j(\mathbf z)\,
 p^Q(\theta^{[j:u]},r)(\mathbf z)\,
 \Lambda_r^Q(d\mathbf z).
\]
The finite polynomial representation and integrable coefficients justify all these integrals, coefficientwise differentiation, and the exchanges of integrals below. In particular,
\[
|\mathsf{sign}^{\theta,r}_j|\le1,
\qquad
\mathsf{sign}^{\theta,r}_j\,\mathsf{diff}^{\theta,r}_j
=|\mathsf{diff}^{\theta,r}_j|,
\]
and
\[
\left(\mathsf{poly}^{\theta,r}_j\right)^{(q)}(u)
=
\int
 \mathsf{sign}^{\theta,r}_j(\mathbf z)\,
 \partial_{\theta_j}^{\,q}p^Q(\theta^{[j:u]},r)(\mathbf z)\,
 \Lambda_r^Q(d\mathbf z).
\]
For \(|u|\le\alpha\), the replacement parameter remains in the cube. Thus
\[
\left|
 \left(\mathsf{poly}^{\theta,r}_j\right)^{(k)}(u)
\right|
\le
\int
 \left|\partial_{\theta_j}^{\,k}
 p^Q(\theta^{[j:u]},r)(\mathbf z)\right|
 \Lambda_r^Q(d\mathbf z)
\le k!\sqrt{\binom nk}\,\beta^k.
\]

Define its Taylor polynomial by
\[
\mathsf{taylor}^{\theta,r}_{j,k}(u)
=
\sum_{q=0}^{k-1}
 \frac{\left(\mathsf{poly}^{\theta,r}_j\right)^{(q)}(0)}{q!}\,u^q,
\qquad u\in\mathbb R.
\]
At \(u=0\), the remainder equals zero. For \(u\ne0\), Taylor's theorem with Lagrange remainder on the segment joining \(0\) and \(u\) gives
\[
\begin{split}
\left|
 \mathsf{poly}^{\theta,r}_j(u)
 -\mathsf{taylor}^{\theta,r}_{j,k}(u)
\right|
&\le
\frac{|u|^k}{k!}
\sup_{u_*\in[\min\{0,u\},\,\max\{0,u\}]}
 \left|
 \left(\mathsf{poly}^{\theta,r}_j\right)^{(k)}(u_*)
 \right|\\
&\le |u|^k\sqrt{\binom nk}\,\beta^k
\qquad(|u|\le\alpha).
\end{split}
\]
Matching moments cancel its Taylor polynomial:
\[
\int\mathsf{taylor}^{\theta,r}_{j,k}(u)\,\nu_0(du)
=
\int\mathsf{taylor}^{\theta,r}_{j,k}(u)\,\nu_1(du).
\]
Each prior is supported on \([-\alpha,\alpha]\) and has mass one, so
\[
\left|
\int
 \left(\mathsf{poly}^{\theta,r}_j(u)
       -\mathsf{taylor}^{\theta,r}_{j,k}(u)\right)\nu_\ell(du)
\right|
\le\alpha^k\sqrt{\binom nk}\,\beta^k,
\qquad \ell\in\{0,1\}.
\]
Finally, exchanging the integrals in the scalar polynomial gives
\[
\begin{split}
\int|\mathsf{diff}^{\theta,r}_j(\mathbf z)|\,\Lambda_r^Q(d\mathbf z)
&=
\int\mathsf{poly}^{\theta,r}_j(u)\,\nu_0(du)
-\int\mathsf{poly}^{\theta,r}_j(u)\,\nu_1(du)\\
&\le 2\alpha^k\sqrt{\binom nk}\,\beta^k.
\end{split}
\]
This is uniform in the remaining coordinates, the seed, and the coordinate being replaced.
% lean: CausalSmith.Stat.LdpOptvalueUniformFrontier.canonical_density_coordinate_matching_L1

\item \textbf{Averaging and coordinate telescoping.}
Define the endpoint product priors and the intervening product priors by
\[
\Pi_\ell=\nu_\ell^{\otimes d},
\qquad \ell\in\{0,1\},
\]
\[
\Pi_v^{\mathrm{hyb}}
=\nu_1^{\otimes v}\otimes\nu_0^{\otimes(d-v)},
\qquad v=0,\ldots,d.
\]
Their coordinates are ordered from \(1\) to \(d\). They are probability laws supported on the parameter cube, with
\[
\Pi_0^{\mathrm{hyb}}=\Pi_0,
\qquad
\Pi_d^{\mathrm{hyb}}=\Pi_1.
\]
For each \(r\) and \(v\), define the averaged density
\[
\mathsf{avg}_v^r(\mathbf z)
=\int p^Q(\theta,r)(\mathbf z)\,\Pi_v^{\mathrm{hyb}}(d\theta).
\]
These functions are integrable against \(\Lambda_r^Q\), by their finite-row representation.

For \(v<d\), set \(j=v+1\). Resampling coordinate \(j\) from \(\nu_0\) preserves \(\Pi_v^{\mathrm{hyb}}\); resampling it from \(\nu_1\) gives \(\Pi_{v+1}^{\mathrm{hyb}}\). Explicitly, with \(\#\) denoting pushforward,
\[
(\theta,u\mapsto\theta^{[j:u]})_\#
 (\Pi_v^{\mathrm{hyb}}\otimes\nu_0)=\Pi_v^{\mathrm{hyb}},
\qquad
(\theta,u\mapsto\theta^{[j:u]})_\#
 (\Pi_v^{\mathrm{hyb}}\otimes\nu_1)=\Pi_{v+1}^{\mathrm{hyb}}.
\]
Fubini therefore gives
\[
\mathsf{avg}_v^r(\mathbf z)-\mathsf{avg}_{v+1}^r(\mathbf z)
=
\int\mathsf{diff}^{\theta,r}_j(\mathbf z)\,
 \Pi_v^{\mathrm{hyb}}(d\theta).
\]
The triangle inequality under the integral, followed by the single-coordinate estimate, yields
\[
\begin{split}
\int
 |\mathsf{avg}_v^r-\mathsf{avg}_{v+1}^r|\,d\Lambda_r^Q
&\le
\int
 \left\{\int|\mathsf{diff}^{\theta,r}_j(\mathbf z)|
              \Lambda_r^Q(d\mathbf z)\right\}
 \Pi_v^{\mathrm{hyb}}(d\theta)\\
&\le2\alpha^k\sqrt{\binom nk}\,\beta^k.
\end{split}
\]
Telescoping the \(d\) replacements now gives
\[
\begin{split}
&\int
 \left|
 \int p^Q(\theta,r)(\mathbf z)\,\Pi_0(d\theta)
 -\int p^Q(\theta,r)(\mathbf z)\,\Pi_1(d\theta)
 \right|
 \Lambda_r^Q(d\mathbf z)\\
&\qquad
=\int|\mathsf{avg}_0^r-\mathsf{avg}_d^r|\,d\Lambda_r^Q\\
&\qquad
\le\sum_{v=0}^{d-1}
 \int|\mathsf{avg}_v^r-\mathsf{avg}_{v+1}^r|\,d\Lambda_r^Q\\
&\qquad
\le2d\alpha^k\sqrt{\binom nk}\,\beta^k.
\end{split}
\]
% lean: CausalSmith.Stat.LdpOptvalueUniformFrontier.density_product_matching_of_gate

\item \textbf{Fixed-seed total variation.}
To define probability kernels on the full parameter domain, set
\[
\theta^{\mathrm{cl}}_j
=\max\left\{-\frac12,\min\left\{\frac12,\theta_j\right\}\right\},
\qquad
\theta\in\mathbb R^d,\ j\in[d].
\]
This parameter belongs to the cube for every \(\theta\), and equals \(\theta\) on the cube. Define, for \(\ell\in\{0,1\}\) and \(r\in\mathcal R\),
\[
\mathsf{condmix}_{\ell,r}(d\mathbf z)
=\int\mathsf{Prob}^Q_{\theta^{\mathrm{cl}},r}(d\mathbf z)\,
 \Pi_\ell(d\theta),
\]
\[
\mathsf{seedrow}_{\ell,r}
=(\mathbf z\mapsto(r,\mathbf z))_\#
 \mathsf{condmix}_{\ell,r},
\qquad
\mathsf{avden}_\ell(r,\mathbf z)
=\int p^Q(\theta,r)(\mathbf z)\,\Pi_\ell(d\theta).
\]
All these measure-valued families are measurable; the first two consist of probability measures. Mixing commutes with the displayed measurable map, so
\[
\mathsf{seedrow}_{\ell,r}
=
\int
 (\mathbf z\mapsto(r,\mathbf z))_\#
 \mathsf{Prob}^Q_{\theta^{\mathrm{cl}},r}\,
 \Pi_\ell(d\theta).
\]
Support on the cube, the density representation, and Fubini give
\[
\mathsf{condmix}_{\ell,r}(d\mathbf z)
=\mathsf{avden}_\ell(r,\mathbf z)\Lambda_r^Q(d\mathbf z),
\qquad
\mathsf{avden}_\ell(r,\mathbf z)\ge0,
\qquad
\int\mathsf{avden}_\ell(r,\mathbf z)\Lambda_r^Q(d\mathbf z)=1.
\]

For completeness, the last normalization gives
\[
\int\bigl(\mathsf{avden}_0(r,\mathbf z)
          -\mathsf{avden}_1(r,\mathbf z)\bigr)
 \Lambda_r^Q(d\mathbf z)=0.
\]
For every measurable \(E\subseteq\mathcal Z\), the integral of this difference over \(E\) is the negative of its integral over \(E^c\). Thus
\[
\begin{split}
&2\left|
 \mathsf{condmix}_{0,r}(E)-\mathsf{condmix}_{1,r}(E)
 \right|\\
&\qquad\le
 \int_E|\mathsf{avden}_0-\mathsf{avden}_1|\,d\Lambda_r^Q
 +\int_{E^c}|\mathsf{avden}_0-\mathsf{avden}_1|\,d\Lambda_r^Q\\
&\qquad=
 \int|\mathsf{avden}_0-\mathsf{avden}_1|\,d\Lambda_r^Q.
\end{split}
\]
Taking the supremum over measurable events gives the half-\(L^1\) bound. A measurable pushforward decreases total variation, since its event differences are differences over measurable preimages. Consequently, at every seed,
\[
\begin{split}
\operatorname{TV}(\mathsf{seedrow}_{0,r},\mathsf{seedrow}_{1,r})
&\le
\operatorname{TV}(\mathsf{condmix}_{0,r},\mathsf{condmix}_{1,r})\\
&\le
\frac12\int
 |\mathsf{avden}_0(r,\mathbf z)-\mathsf{avden}_1(r,\mathbf z)|
 \Lambda_r^Q(d\mathbf z)\\
&\le d\alpha^k\sqrt{\binom nk}\,\beta^k.
\end{split}
\]
% lean: CausalSmith.Stat.LdpOptvalueUniformFrontier.conditionalMixture_tv_le_half_L1

\item \textbf{Mixing over the common public seed.}
The sampling and randomness assumptions ensure that, at each cube parameter, the observed-record vector has its independent product law and is independent of the public seed and analyst randomizer. Successive kernel integration therefore gives the seed-transcript law
\[
\lambda(dr)\,\mathsf{Prob}^Q_{\theta,r}(d\mathbf z).
\]
The priors are supported on the cube and are fixed independently of the seed. Hence both \(\mathsf M_\ell^Q\) are probability measures, and
\[
\operatorname{TV}(\mathsf M_0^Q,\mathsf M_1^Q)\le1.
\]
Interchanging the independent prior and seed integrations gives, for every measurable seed-transcript event \(E\),
\[
\begin{split}
\mathsf M_\ell^Q(E)
&=
\int_{\mathbb R^d}\int_{\mathcal R}
 \left[
 (\mathbf z\mapsto(r,\mathbf z))_\#
 \mathsf{Prob}^Q_{\theta^{\mathrm{cl}},r}
 \right](E)\,\lambda(dr)\,\Pi_\ell(d\theta)\\
&=
\int_{\mathcal R}\mathsf{seedrow}_{\ell,r}(E)\,\lambda(dr).
\end{split}
\]
The measurability established above justifies these mixture identities. Applying the fixed-seed bound to each event gives
\[
\begin{split}
|\mathsf M_0^Q(E)-\mathsf M_1^Q(E)|
&\le
\int
 |\mathsf{seedrow}_{0,r}(E)-\mathsf{seedrow}_{1,r}(E)|
 \,\lambda(dr)\\
&\le
\int d\alpha^k\sqrt{\binom nk}\,\beta^k\,\lambda(dr)\\
&=d\alpha^k\sqrt{\binom nk}\,\beta^k.
\end{split}
\]
Taking the supremum and combining with the probability-measure bound proves
\[
\operatorname{TV}(\mathsf M_0^Q,\mathsf M_1^Q)
\le
\min\left\{1,\,
d\alpha^k\sqrt{\binom nk}\,\beta^k\right\},
\qquad 1\le k\le n.
\]
% lean: CausalSmith.Stat.LdpOptvalueUniformFrontier.probability_bind_tv_le_constant

\item \textbf{Exact cancellation above the sample degree.}
Now suppose \(k>n\). Define the prior-predictive input laws on
\(\mathcal O_d^n\) by
\[
\mathsf{Input}_\ell
=
\int
 \left(\sum_{o\in\mathcal O_d}
  \mathsf{mass}^{(1)}_{\theta^{\mathrm{cl}}}(o)\,\delta_o
 \right)^{\otimes n}
 \Pi_\ell(d\theta),
\qquad \ell\in\{0,1\}.
\]
Since clipping equals the identity on the prior support, their atom masses are
\[
\mathsf{Input}_\ell(\{\mathbf o\})
=\int\mathsf{mass}^{(n)}_\theta(\mathbf o)\,\Pi_\ell(d\theta).
\]
For \(\mathcal J\subseteq[n]\), \(\mathbf o\in\mathcal O_d^n\), and \(j\in[d]\), define the coordinate counts
\[
\mathsf{count}_{j,\mathcal J}(\mathbf o)
=\#\{i\in\mathcal J:j_{o_i}=j\}.
\]
Expanding the affine sample mass gives
\[
\mathsf{mass}^{(n)}_\theta(\mathbf o)
=
(4d)^{-n}
\sum_{\mathcal J\subseteq[n]}
 \left\{\prod_{i\in\mathcal J}\mathsf{sgn}_{\mathrm{obs}}(o_i)\right\}
 \prod_{i\in\mathcal J}\theta_{j_{o_i}}.
\]
Regrouping factors by coordinate and integrating under the product prior yields
\[
\begin{split}
\int\prod_{i\in\mathcal J}\theta_{j_{o_i}}\,\Pi_\ell(d\theta)
&=
\int\prod_{j=1}^d
 \theta_j^{\mathsf{count}_{j,\mathcal J}(\mathbf o)}
 \,\Pi_\ell(d\theta)\\
&=
\prod_{j=1}^d
 \int u^{\mathsf{count}_{j,\mathcal J}(\mathbf o)}\,\nu_\ell(du).
\end{split}
\]
Every count lies between \(0\) and \(n<k\), so every factor is the same for the two priors by moment matching. This also covers the empty subset and zero counts. Finite expansion therefore gives
\[
\mathsf{Input}_0(\{\mathbf o\})
=\mathsf{Input}_1(\{\mathbf o\})
\qquad
\text{for every \(\mathbf o\in\mathcal O_d^n\)},
\]
and hence
\[
\mathsf{Input}_0=\mathsf{Input}_1.
\]

Define the probability channel from a fixed input vector to the seed-transcript space by
\[
\mathsf{channel}^Q_{\mathbf o}(dr,d\mathbf z)
=\lambda(dr)\,\mathsf{Law}^Q_{\mathbf o,r}(d\mathbf z).
\]
The sampling and randomness assumptions give
\[
\mathsf M_\ell^Q
=
\sum_{\mathbf o\in\mathcal O_d^n}
 \mathsf{Input}_\ell(\{\mathbf o\})\,
 \mathsf{channel}^Q_{\mathbf o}.
\]
The independent analyst randomizer integrates to mass one in this representation. Applying the same channel to the identical input mixtures proves
\[
\mathsf M_0^Q=\mathsf M_1^Q
\qquad(k>n).
\]
% lean: CausalSmith.Stat.LdpOptvalueUniformFrontier.mixtureLaw_eq_of_matching_above_sample

\item \textbf{Bounded transcript decisions.}
Both mixtures are probability measures for every supported pair of priors. The standard bounded-range total-variation inequality, applied to a Borel measurable
\(f:\mathcal R\times\mathcal Z\to[0,1]\), gives
\[
\left|
\int f(r,\mathbf z)\,\mathsf M_0^Q(dr,d\mathbf z)
-
\int f(r,\mathbf z)\,\mathsf M_1^Q(dr,d\mathbf z)
\right|
\le
\operatorname{TV}(\mathsf M_0^Q,\mathsf M_1^Q).
\]
This applies for every \(k\ge1\) in the statement.
% lean: CausalSmith.Stat.LdpOptvalueUniformFrontier.mixtureLaw_bounded_test
\end{enumerate}
\end{proof}
